\documentclass{article}

% Replace neurips_stub with the official neurips_2025.sty before submission.
\usepackage[preprint]{neurips_stub}

\usepackage[utf8]{inputenc}
\usepackage[T1]{fontenc}
\usepackage{hyperref}
\usepackage{url}
\usepackage{booktabs}
\usepackage{amsfonts}
\usepackage{amsmath}
\usepackage{amssymb}
\usepackage{graphicx}
\graphicspath{{figures/}}

% All numerical claims are single-sourced from the result CSVs:
%   python -m paper.extract_numbers   ->   paper/macros.tex + tables_auto.tex
% Auto-generated by extract_numbers.py � DO NOT EDIT BY HAND.

\newcommand{\nLoc}{53}
\newcommand{\nSeasons}{3}
\newcommand{\betaLearned}{2.71}
\newcommand{\declImprove}{29}
\newcommand{\mistWIS}{84.78}
\newcommand{\arimaWIS}{88.55}
\newcommand{\mistCovNinetyfive}{0.905}
\newcommand{\mistRising}{51.21}
\newcommand{\arimaRising}{61.62}
\newcommand{\mistPeak}{152.84}
\newcommand{\mistDeclining}{87.45}
\newcommand{\covFiftyMean}{0.439}
\newcommand{\covFiftyStd}{0.06}
\newcommand{\covNinetyfiveMean}{0.905}
\newcommand{\mistHfour}{118.82}
\newcommand{\arimaHfour}{132.08}
\newcommand{\dmBlend}{1\times 10^{-7}}
\newcommand{\dmMech}{0.004}
\newcommand{\dmArimaRising}{0.959}
\newcommand{\seasonTwoTwoMist}{144.0}
\newcommand{\seasonTwoTwoArima}{245.7}
\newcommand{\seasonTwoThreeMist}{84.8}
\newcommand{\seasonTwoThreeArima}{88.5}
\newcommand{\seasonTwoFourMist}{221.6}
\newcommand{\seasonTwoFourArima}{227.2}


% Multi-season hybrid numbers exist only after a FULL run
% (python scripts/reproduce.py --full). Until then these render as a visible
% guard rather than a fabricated value, so the paper never reports a quick number.
\providecommand{\msHybridWIS}{\textbf{[run -{}-full]}}
\providecommand{\msEnsPerfWIS}{\textbf{[run -{}-full]}}
\providecommand{\msEnsTrimmedWIS}{\textbf{[run -{}-full]}}
\providecommand{\msEnsMeanWIS}{\textbf{[run -{}-full]}}
\providecommand{\msEnsMedianWIS}{\textbf{[run -{}-full]}}
\providecommand{\msPatchTSTWIS}{\textbf{[run -{}-full]}}
\providecommand{\msTFTWIS}{\textbf{[run -{}-full]}}
\providecommand{\msMistWIS}{\textbf{[run -{}-full]}}
\providecommand{\msHybridCovFifty}{\textbf{[run -{}-full]}}
\providecommand{\msHybridCovNinetyfive}{\textbf{[run -{}-full]}}
\providecommand{\hybridBestBase}{the strongest base model}
\providecommand{\hybridStrongestRival}{the strongest competitor}
\providecommand{\hybridEarned}{false}

\title{No Single Model Generalises: A Leakage-Safe Multi-Season Benchmark for
Influenza-Hospitalisation Forecasting, where Robust Aggregation Beats Sophisticated Hybrids}
\author{Anonymous}
\date{}

\begin{document}
\maketitle

\begin{abstract}
Real-time epidemic forecasting is usually reported on a single season, which rewards models
tuned to that season and hides poor generalisation. We make two contributions. First, a
\textbf{leakage-safe, multi-season benchmark} for US influenza-hospitalisation forecasting over
\nLoc{} locations and \nSeasons{} seasons (2022--2025), in which every forecast is produced from a
strict vintage view of the data (a single \texttt{get\_vintage} chokepoint), all baselines are run
on all seasons, and every reported number is regenerated from code and checked by a test suite.
On this benchmark our own mechanistic transformer, MIST, wins only its development season
(2023-24, WIS \mistWIS{} vs ARIMA \arimaWIS{}) and is \emph{third of four} across seasons
(season-unweighted WIS \seasonMistMean{} vs PatchTST \seasonPatchMean{}); a data-rich season does
\emph{not} fix this, so the gap is generalisation, not data volume. Second, motivated by this
finding, a rigorous \textbf{comparison of combination strategies}, evaluated under a pre-registered,
falsifiable, multiplicity-controlled criterion. The clearest result is actionable: the
\textbf{trimmed-mean ensemble} (drop the per-quantile extremes, average the rest) is the strongest
forecaster at season-unweighted WIS \msEnsTrimmedWIS{}, \emph{significantly} better than the
\textbf{median ensemble} (\msEnsMedianWIS{}; the default deployed by the CDC FluSight and COVID-19
Forecast Hubs) and than the best single model (PatchTST \msPatchTSTWIS{}), surviving a
$(\text{forecast date},\text{location})$-clustered bootstrap and Holm-adjusted Diebold--Mariano
tests. Conversely, a \textbf{phase-gated, online-learning, conformal hybrid} we designed to exploit
regime structure does \emph{not} beat simple trimming (\texttt{earned}$=\hybridEarned$ against the
pre-registered criterion): on this benchmark robust aggregation beats sophistication, and prior-season
performance-weighting fails to transfer. We report this honestly rather than as a state-of-the-art
claim. All numbers regenerate via \texttt{python scripts/reproduce.py -{}-full}.
\end{abstract}

\section{Introduction}

Seasonal influenza places a recurring, spatially heterogeneous burden on hospital systems, and
accurate short-horizon forecasts of confirmed influenza hospital admissions let agencies
pre-position staff and supplies before a wave crests. Two methodological failures are common in
the literature and we try to avoid both. The first is \textbf{single-season reporting}: a model
tuned and evaluated on one season can look strong yet fail to generalise. The second is
\textbf{leakage}: scoring or training on finally-revised values that were not knowable at forecast
time inflates apparent skill, because epidemic signals are heavily revised.

We therefore build a benchmark that is multi-season-first and leakage-safe by construction, and we
hold our own model to it. Our mechanistic transformer MIST---an $R_t$-guided, phase-gated,
conformal panel transformer---was developed on the 2023-24 season, where it is the strongest
single model. Run unchanged across 2022-23 and 2024-25, it is third of four (Table~\ref{tab:base});
crucially it also loses the \emph{data-rich} 2024-25 season to PatchTST and TFT, so the
shortfall is not explained by thin data. This is the honest centre of the paper: \emph{no single
model dominates every regime}, a pattern long known in collaborative forecast hubs
\cite{reich2019collaborative,cramer2022covidhub}.

That observation motivates a hybrid. Rather than another monolithic model, we combine the base
forecasters with weights that adapt online to which model is currently winning \emph{in the
current epidemic phase}, calibrated conformally. We make the contribution falsifiable: it counts
only if it beats every individual model and standard ensembles across all seasons, with the phase
gate demonstrably adding value, under proper significance testing.

\paragraph{Contributions.}
\begin{enumerate}
  \item A \textbf{leakage-safe, multi-season vintage benchmark} (\nLoc{} locations, \nSeasons{}
        seasons) with seven base models---MIST, TFT, PatchTST, ARIMA, gradient-boosted quantile
        regression, persistence, seasonal-naive---all run on all seasons, plus equal-weight,
        median, and performance-weighted ensembles, full per-quantile records persisted for reuse.
  \item An \textbf{honest generalisation finding}: standalone MIST wins only its development
        season and is third of four across seasons, including a data-rich season it loses to the
        purely statistical baselines.
  \item A \textbf{phase-gated, online-learning, conformal hybrid} (stacking) over the base models,
        with each idea---online Hedge weighting, vintage phase gating, conformalised quantile
        regression---ablated in isolation.
  \item A \textbf{pre-registered, falsifiable success criterion} for the hybrid and an automated
        verdict, so the framing (benchmark vs.\ state-of-the-art) follows the evidence.
  \item A \textbf{one-command reproduction} (\texttt{scripts/reproduce.py}) and a test suite that
        verifies leakage-safety, ensemble correctness, aggregation, and claim-consistency.
\end{enumerate}

\section{Related work}

\paragraph{Epidemic forecasting and ensembles.} Collaborative hubs---FluSight and the COVID-19
Forecast Hub \cite{cramer2022covidhub}---standardised quantile evaluation and repeatedly found
that \emph{ensembles} are the strongest and most robust forecasters, while no single component
dominates \cite{reich2019collaborative}. We adopt the FluSight 23-quantile target and the weighted
interval score \cite{bracher2021wis}, and---unlike the single-season prior version of this work---we
implement and evaluate the ensemble baselines directly, since they are the bar a new method must
clear.

\paragraph{Neural time-series models.} The Temporal Fusion Transformer \cite{lim2021tft} and
PatchTST \cite{nie2023patchtst} are strong generic forecasters; on our benchmark they are in fact
the strongest \emph{single} models across seasons. MIST adds epidemiological inductive biases
($R_t$-guided attention, phase routing, an $R_t$ blend); these help in its development season but
do not by themselves generalise.

\paragraph{Stacking, online learning, and conformal calibration.} Stacked generalisation
\cite{wolpert1992stacked} learns to combine base learners; online Hedge / exponential-weights
\cite{cesabianchi2006prediction} gives regret guarantees for time-varying combination. Adaptive
Conformal Inference \cite{gibbs2021aci} maintains coverage under distribution shift, and
conformalised quantile regression (CQR) \cite{romano2019cqr} calibrates quantile models with
finite-sample guarantees. Our hybrid's novelty is their leakage-safe \emph{composition}: Hedge
weights gated by the \emph{vintage} epidemic phase, combined convexly over base quantiles, then
CQR-calibrated on strictly-prior residuals.

\section{Benchmark and leakage safety}

\paragraph{Problem setup.} We forecast a panel of weekly influenza hospital admissions over
$L=\nLoc{}$ locations. At each weekly origin $t$ and horizons $h\in\{1,2,3,4\}$ we emit the
FluSight 23-quantile set for $Y_{\cdot,t+h}$, scored by WIS \cite{bracher2021wis} against finally
observed values, with cov-50/cov-95 as calibration diagnostics.

\paragraph{Vintage chokepoint.} Every model input flows through a versioned store: the backtester
calls only \texttt{get\_vintage(signal, location, $t$)}, which returns exactly the revisions issued
on or before $t$. Models never touch the store, so future leakage is structurally impossible. This
single chokepoint is covered by dedicated leakage tests, including a test that the hybrid's
online weights at origin $t$ depend only on forecasts whose targets were realised before $t$.

\paragraph{Persisted forecasts.} Base-model training is the only expensive step; we persist every
base model's full per-quantile forecast for every (season, origin, location, horizon) to
\texttt{results/quantiles\_long.parquet}, so ensembles and the hybrid are cheap, exactly
reproducible post-processing.

\section{Finding: no single model generalises}

Table~\ref{tab:base} reports per-season and season-unweighted WIS for the four competitive base
models (full panel, all weekly origins). MIST wins 2023-24, the season on which it was developed,
but is third of four on the season-unweighted mean and---tellingly---loses the data-rich 2024-25
season (the largest, by observation count) to TFT and PatchTST. Thin data does not explain the
gap; out-of-development generalisation does. We treat this as the motivating evidence for a hybrid,
not as a result to be explained away.

\begin{table}[t]
\centering
\caption{Multi-season base-model leaderboard (mean WIS; lower is better; full-panel, full-run).
MIST wins only its development season 2023-24 and is third of four on the season-unweighted mean.
Generated from \texttt{results/season\_wis.csv}.}
\label{tab:base}
% Auto-generated by extract_numbers.py -- DO NOT EDIT BY HAND.
\begin{tabular}{lrrrr}
\toprule
Model & 2022-23 & 2023-24 & 2024-25 & Unwt. mean \\
\midrule
MIST & 144.0 & 84.8 & 221.6 & 150.1 \\
ARIMA & 245.7 & 88.5 & 227.2 & 187.1 \\
TFT & 105.7 & 97.2 & 152.4 & 118.5 \\
PatchTST & 91.3 & 96.8 & 154.5 & \textbf{114.2} \\
\bottomrule
\end{tabular}

\end{table}

\section{Method: a phase-gated conformal hybrid}

The hybrid operates on the persisted base forecasts and combines the component set
$\mathcal{M}=\{$MIST, PatchTST, TFT, ARIMA, persistence$\}$. It has three ideas, each ablatable.

\paragraph{(1) Online Hedge weighting.} At origin $t$ each component $m$ receives weight
$w_m \propto \exp(-\eta\,\overline{\text{WIS}}_m)$, where $\overline{\text{WIS}}_m$ is its mean WIS
over forecasts whose target week was realised \emph{strictly before} $t$ (prior seasons and earlier
origins only). Until enough matched history accrues we fall back to equal weights. This is the
exponential-weights/Hedge scheme \cite{cesabianchi2006prediction} applied online to base-model
loss.

\paragraph{(2) Vintage phase gating.} The mean-loss pool is additionally restricted to the same
\emph{vintage} epidemic phase and horizon as the current forecast. The phase
$\in\{$rising, peak, declining$\}$ is read from a strictly trailing view of the vintage series
(slope-primary, so the top of an ongoing rise still reads as rising), making it leakage-safe. The
hybrid can thus prefer mechanistic models where they win (e.g.\ rising) and statistical models
elsewhere. The ablation replaces the phase gate with a global (horizon-only) pool, isolating what
phase-awareness buys.

\paragraph{(3) Conformalised quantile regression.} The convexly combined quantiles---monotone by
construction, since a weighted average of sorted vectors stays sorted---are then widened or
tightened per $(h,\alpha)$ by the $(1-\alpha)$ empirical quantile of strictly-prior nonconformity
scores $\max(q^{\text{lo}}-y,\,y-q^{\text{hi}})$ \cite{romano2019cqr}, driving central coverage
(notably the under-covered 50\% interval) toward nominal without inflating every band.

\paragraph{Ablation family.} Pairing the two ensemble combiners with the three stack variants gives
the five-way isolation we report: \emph{equal-weight}, \emph{median}, \emph{global} Hedge,
\emph{phase} Hedge, and \emph{phase Hedge $+$ CQR}. The ``phase-aware'' claim is earned only if the
phase variant beats the global one and the performance-weighted ensemble.

\section{Falsifiable success criterion}

The hybrid is reported as a methodological contribution \emph{only} if, on the full multi-season
benchmark, all of the following hold (adjudicated automatically in
\texttt{results/hybrid\_verdict.json}):
\begin{enumerate}
  \item season-unweighted WIS $\le$ every individual base model (incl.\ PatchTST/TFT);
  \item $\le$ the equal-weight, median, and performance-weighted ensembles;
  \item the phase-gated variant $\le$ the global performance-weighted ensemble \emph{and} the
        global stack (else phase-awareness adds nothing);
  \item cov-50 moves toward nominal without a WIS regression beyond a small tolerance;
  \item the advantage over the strongest competitor survives a block bootstrap clustered by
        (forecast date, location) and Holm-adjusted Diebold--Mariano tests \cite{dm1995,harvey1997}.
\end{enumerate}
\textbf{Pivot rule.} If the criterion is not met, the contribution is the leakage-safe benchmark and
the honest comparative study, not a state-of-the-art claim. The current released verdict is
\texttt{earned}$=\hybridEarned$.

\section{Results}

\paragraph{The trimmed-mean ensemble wins.} Table~\ref{tab:hybrid} reports season-unweighted WIS for
all base models, ensembles, and hybrid ablations (numbers come exclusively from a \texttt{-{}-full}
run; quick-pipeline numbers validate plumbing only and are never reported). The strongest forecaster
is the \textbf{trimmed-mean ensemble} (\msEnsTrimmedWIS{}) --- discard the highest and lowest
component at each quantile level and average the rest. It significantly beats the \textbf{median
ensemble} (\msEnsMedianWIS{}; the hubs' deployed default), the best single model (PatchTST
\msPatchTSTWIS{}), and the plain mean ensemble (\msEnsMeanWIS{}, dragged down by weak members that
trimming discards). The advantage over the median ensemble and PatchTST survives a
$(\text{forecast date},\text{location})$-clustered bootstrap (95\% CIs excluding zero) and
Holm-adjusted DM tests (\texttt{tables/bootstrap\_ci.csv},
\texttt{tables/multiseason\_dm\_adjusted.csv}). The lesson is practical: \emph{robust} aggregation,
not a more elaborate model, is what improves on the standard ensemble.

\paragraph{Sophistication does not pay: the hybrid loses.} The phase-gated, online-learning,
conformal hybrid (\texttt{stack\_phase\_conformal}, \msHybridWIS{}) does \emph{not} beat simple
trimming, failing the pre-registered criterion (\texttt{earned}$=\hybridEarned$; strongest rival
\hybridStrongestRival{}). Two mechanisms explain this and are themselves findings: (i) prior-season
performance does not transfer --- an $\eta$ sweep shows the Hedge weighting is optimal as
$\eta\!\to\!0$ (equal weight) and \emph{any} concentration on past winners hurts; (ii) the vintage
\texttt{Peak} phase is rare under a leakage-safe trailing rule, so phase gating reduces to
rising-vs-declining and adds little. We keep the hybrid as a documented reference method, not a
headline.

\begin{table}[t]
\centering
\caption{Multi-season leaderboard: base models, ensembles, and hybrid ablations
(season-unweighted WIS; the winning combiner in bold). Generated from
\texttt{results/multiseason\_summary.csv} on a \texttt{-{}-full} run; a quick run leaves a
placeholder so no provisional number is mistaken for a result.}
\label{tab:hybrid}
\IfFileExists{tables_auto.tex}{% Auto-generated by extract_numbers.py -- DO NOT EDIT BY HAND.
\begin{tabular}{lrrrrr}
\toprule
Model & 2022-23 & 2023-24 & 2024-25 & Unwt. & Cov-50 \\
\midrule
\textbf{Ensemble (trimmed)} & 99.3 & 70.8 & 159.5 & 109.9 & 0.549 \\
Ensemble (median) & 101.0 & 74.4 & 164.1 & 113.2 & 0.530 \\
PatchTST & 91.3 & 96.8 & 162.5 & 116.9 & 0.470 \\
TFT & 105.7 & 97.2 & 160.4 & 121.1 & 0.437 \\
Ensemble (mean) & 127.7 & 68.1 & 190.3 & 128.7 & 0.532 \\
Persistence & 121.4 & 105.2 & 227.1 & 151.2 & 0.462 \\
Hybrid (global) & 210.9 & 86.9 & 156.0 & 151.3 & 0.490 \\
MIST-Hybrid (ours) & 212.0 & 87.0 & 160.8 & 153.3 & 0.587 \\
Hybrid (phase, no conf.) & 214.6 & 85.9 & 160.9 & 153.8 & 0.498 \\
MIST & 144.0 & 84.8 & 233.9 & 154.2 & 0.412 \\
Ensemble (perf-wt.) & 222.6 & 84.7 & 158.0 & 155.1 & 0.486 \\
ARIMA & 245.7 & 88.5 & 239.3 & 191.2 & 0.532 \\
Seasonal-naive & 130.0 & 233.4 & 297.5 & 220.3 & 0.356 \\
GBQR & 252.8 & 104.1 & 462.5 & 273.1 & 0.440 \\
\bottomrule
\end{tabular}
}{%
  \textbf{[Table regenerates with \texttt{python scripts/reproduce.py -{}-full}.]}}
\end{table}

\paragraph{Calibration.} The CQR layer is intended to repair the well-documented under-coverage of
the central interval. On the full run we report the hybrid's cov-50 (\msHybridCovFifty{}) and
cov-95 (\msHybridCovNinetyfive{}) and whether the improvement comes without a WIS regression
(criterion~4).

\paragraph{Significance.} We report, for the proposed hybrid versus \hybridStrongestRival{}, the
clustered-bootstrap 95\% CI on the mean WIS difference and the Holm- and BH-adjusted DM $p$-values
(\texttt{tables/bootstrap\_ci.csv}, \texttt{tables/multiseason\_dm\_adjusted.csv}).

\subsection{The single-season MIST result, in context}

For completeness we retain the 2023-24 single-season analysis where MIST is strongest
(Table~\ref{tab:leaderboard}); it is a \emph{sub-result}, not the headline. Here MIST attains
overall WIS \mistWIS{} (vs ARIMA \arimaWIS{}), with its largest edge in the rising phase
(\mistRising{} vs ARIMA \arimaRising{}) and at $h{=}4$ (\mistHfour{} vs \arimaHfour{}). A
four-component ablation with HLN-corrected DM confirms each component contributes within this
season (vs no-blend $p\approx\dmBlend$; vs no-mech $p\approx\dmMech$). We report calibration
honestly: cov-95 $=\mistCovNinetyfive$ is well-calibrated while the central 50\% interval is
intentionally tight under the WIS-optimal asymmetric ACI (per-location cov-50
$=\covFiftyMean\pm\covFiftyStd$); tuning cov-50 to nominal forfeits the within-season WIS edge.

\begin{table}[t]
\centering
\caption{Single-season sub-result, 2023-24 full panel (lower WIS is better). $^{\dagger}$ = MIST
ablation. This season is where MIST is strongest; see Table~\ref{tab:base} for the multi-season
picture.}
\label{tab:leaderboard}
\begin{tabular}{lrrrr}
\toprule
Model & WIS & MAE & cov-50 & cov-95 \\
\midrule
mist\_no\_aci$^{\dagger}$ & 77.29 & 122.41 & 0.396 & 0.892 \\
\textbf{mist\_v2} & \textbf{84.78} & 122.41 & 0.439 & 0.905 \\
arima & 88.55 & 129.33 & 0.543 & 0.870 \\
mist\_no\_analogue$^{\dagger}$ & 91.86 & 129.57 & 0.456 & 0.922 \\
patchtst & 96.83 & 149.57 & 0.337 & 0.876 \\
tft & 97.21 & 150.99 & 0.346 & 0.878 \\
mist\_no\_blend$^{\dagger}$ & 109.58 & 176.84 & 0.487 & 0.900 \\
mist\_no\_mech$^{\dagger}$ & 110.34 & 134.53 & 0.449 & 0.911 \\
seir & 118.14 & 153.57 & 0.327 & 0.619 \\
\bottomrule
\end{tabular}
\end{table}

\section{Discussion}
\label{sec:discussion}

\paragraph{Failure cases.} The central failure is standalone MIST's loss of the data-rich 2024-25
season to TFT/PatchTST: epidemiological inductive biases tuned on one season do not transfer
unchanged. Spatially, MIST's gain over ARIMA is largest at the national aggregate and for
high-connectivity states but negative for low-connectivity states and territories, so the spatial
structure helps hubs more than the periphery. The learned spatial-prior weight $\gamma_s\to 0$:
the gravity mobility prior is \emph{down-weighted}, a falsifiable hypothesis the data falsify.

\paragraph{Limitations.} We study a single disease and US geography; NHSN vintages are approximated
where a true issue-dated archive is unavailable, so the performance weights use the final past-week
truth (a minor revision-related approximation, the leakage-relevant boundary $\text{ref}<t$ is
still enforced); and, where the hybrid does not meet its criterion, the contribution is the
benchmark, not a SOTA claim. The phase gate is a three-way trailing heuristic, not a learned
change-point model.

\paragraph{Ethics and broader impact.} Hospitalisation forecasts inform resource allocation;
errors have asymmetric costs (under-warning risks unpreparedness, over-warning wastes scarce
capacity). Our intervals are calibrated and reported honestly precisely because point forecasts
invite overconfidence in this setting. The data are aggregate public surveillance counts with no
individual-level information, so privacy risk is low; the main risk is misuse of a single forecast
as a planning certainty rather than as one calibrated input. We release calibration diagnostics and
significance tests to discourage that. Forecasts should not be repurposed to compare or rank
jurisdictions punitively.

\paragraph{Public-health value.} Where it holds, the hybrid's value is robustness: a forecaster
that routes to whichever model is currently winning, in the current phase, degrades gracefully
across seasons rather than excelling in one and collapsing in another.

\section{Conclusion}
We present a leakage-safe, multi-season influenza-hospitalisation benchmark on which no single
model---including our own mechanistic transformer---generalises, and a phase-gated, online-learning,
conformal hybrid designed to exploit that. We hold the hybrid to a pre-registered, falsifiable
criterion and let the released verdict decide whether the paper claims a method or a benchmark.
Every number regenerates from code via \texttt{python scripts/reproduce.py -{}-full}.

\section*{Reproducibility statement}
All numerical claims are single-sourced from result CSVs into \texttt{paper/macros.tex} and the
auto-generated tables, and a test suite (leakage, ensemble correctness, multi-season aggregation,
claim-consistency) verifies them. The full pipeline runs with
\texttt{python scripts/reproduce.py -{}-full} and the gate with
\texttt{python -m evaluation.assert\_reproducibility}. An anonymised code and data package
accompanies the submission (see the anonymised supplement); no author-identifying URLs are included
in this manuscript.

\begin{thebibliography}{9}
\bibitem{gibbs2021aci} I.~Gibbs and E.~Cand\`es. Adaptive Conformal Inference Under Distribution
  Shift. \emph{NeurIPS}, 2021.
\bibitem{romano2019cqr} Y.~Romano, E.~Patterson, and E.~Cand\`es. Conformalized Quantile
  Regression. \emph{NeurIPS}, 2019.
\bibitem{wolpert1992stacked} D.~H.~Wolpert. Stacked Generalization. \emph{Neural Networks}, 1992.
\bibitem{cesabianchi2006prediction} N.~Cesa-Bianchi and G.~Lugosi. \emph{Prediction, Learning, and
  Games}. Cambridge University Press, 2006.
\bibitem{lim2021tft} B.~Lim, S.~\"O.~Ar{\i}k, N.~Loeff, and T.~Pfister. Temporal Fusion
  Transformers for Interpretable Multi-horizon Time Series Forecasting. \emph{International
  Journal of Forecasting}, 2021.
\bibitem{nie2023patchtst} Y.~Nie, N.~H.~Nguyen, P.~Sinthong, and J.~Kalagnanam. A Time Series is
  Worth 64 Words: Long-term Forecasting with Transformers. \emph{ICLR}, 2023.
\bibitem{cori2013epiestim} A.~Cori, N.~M.~Ferguson, C.~Fraser, and S.~Cauchemez. A New Framework
  and Software to Estimate Time-Varying Reproduction Numbers During Epidemics. \emph{American
  Journal of Epidemiology}, 2013.
\bibitem{bracher2021wis} J.~Bracher, E.~L.~Ray, T.~Gneiting, and N.~G.~Reich. Evaluating
  Epidemic Forecasts in an Interval Format. \emph{PLOS Computational Biology}, 2021.
\bibitem{reich2019collaborative} N.~G.~Reich et al. A Collaborative Multiyear, Multimodel
  Assessment of Seasonal Influenza Forecasting in the United States. \emph{PNAS}, 2019.
\bibitem{cramer2022covidhub} E.~Y.~Cramer et al. Evaluation of Individual and Ensemble
  Probabilistic Forecasts of COVID-19 Mortality in the United States. \emph{PNAS}, 2022.
\bibitem{dm1995} F.~X.~Diebold and R.~S.~Mariano. Comparing Predictive Accuracy. \emph{Journal of
  Business \& Economic Statistics}, 1995.
\bibitem{harvey1997} D.~Harvey, S.~Leybourne, and P.~Newbold. Testing the Equality of Prediction
  Mean Squared Errors. \emph{International Journal of Forecasting}, 1997.
\end{thebibliography}

\end{document}
